\sectionlink{序論}
有効場理論のアプローチ，すなわち基礎となる対称性と両立する最低次元の演算子によって系を記述する方法は，素粒子物理学から物性物理学まで，多くの分野で大きな成果を上げてきた．本論文の目的は，この方法を適用し，インフレーションを起こしている宇宙論的背景のまわりの揺らぎの理論を記述することである．

単一場インフレーション模型を調べる通常の方法では，スカラー場$\phi$のラグランジアンから出発し，$\phi$の運動方程式とFRW計量に対するフリードマン方程式を連立して解く．ここで関心があるのは，スカラー場が一様で時間依存する解$\phi_0(t)$に従い，ハッブル・パラメータがゆっくり変化する加速膨張，すなわちインフレーション解である．そのうえで，この背景解のまわりの摂動を調べ，さまざまな宇宙論的観測量に対する予言を求める．

時間発展する解のまわりの摂動の理論は，出発点とした$\phi$の理論とは大きく異なる．$\phi$はすべての微分同相変換（diffs）のもとでスカラーであるのに対し，摂動$\delta\phi$は空間微分同相変換のもとでのみスカラーであり，時間微分同相変換に対しては非線形に変換する：
\be
t \to t + \xi^0(t, \vec x)   \qquad \delta\phi \to \delta\phi + \dot\phi_0(t) \xi^0 \;.
\ee
とくに，$\phi(t,\vec x)=\phi_0(t)$というゲージを選べば，インフラトンの摂動はなくなり，すべての自由度が計量に含まれる．スカラー変数$\delta\phi$は重力子に吸収され，重力子はスカラー・モードと2つのテンソル・ヘリシティという3つの自由度を持つようになる．この現象は，自発的に破れたゲージ理論で起こることと類似している．ゲージ対称性のもとで非線形に変換するゴールドストーン・モードは，ゲージ・ボソンに吸収され（ユニタリーゲージ），質量を持つスピン1粒子を与える．ゴールドストーン場の非線形シグマ模型は，標準模型のヒッグス・セクターのように，ゲージ対称性の線形表現へ埋め込んで紫外（UV）完備化することができる．これは，微分同相変換を線形に表現する$\phi$のラグランジアンから出発する，通常のインフレーションの定式化に対応する．本論文では，時間微分同相変換が非線形に実現される時間発展する真空のまわりで，インフレーション中の摂動の理論を直接記述するという，もう一つの見方を強調したい．この形式は，ヌル・エネルギー条件を破る可能性を調べるために，文献\cite{Creminelli:2006xe}で一般のFRW背景に対して初めて導入された．ここではインフレーション解に焦点を当てて，この形式を拡張する．

ハッブル率$H(t)$を持つ平坦なFRW背景のまわりの摂動を記述し，時間微分同相変換が破れ，空間微分同相変換は破れていない最も一般的なラグランジアンは，ユニタリーゲージで次のように与えられることを示す：
\begin{eqnarray}
\label{eq:Laguni}
S & = & \int  \! d^4 x  \: \sqrt{- g} \;\Big[\frac12 M_{\rm Pl}^2 R+
M_{\rm Pl}^2 \dot{H} g^{00} -  M_{\rm Pl}^2 \left(3
H^2+\dot{H}\right) + \frac{M_2(t)^4}{2!} (g^{00}+1)^2 \\ & & \nonumber + \frac{M_3(t)^4}{3!} (g^{00}+1)^3 + \ldots
-\frac{\bar{M_2}(t)^2}{2}\delta K^\mu {}_\mu {}^2 +... \Big] \; .
\end{eqnarray}
アインシュタイン・ヒルベルト項の次にある2つの演算子は，指定された非摂動解$H(t)$を持つという要請によって固定される．一方，それ以外の演算子はすべて自由であり，同じ背景解を持つさまざまな摂動の理論をすべてパラメータ化する．時間微分同相変換が破れているので，たとえば$g^{00}$や時間一定面の外的曲率$K^\mu {}_\nu {}$など，空間微分同相変換を尊重する任意の項を書くことができる．演算子の係数は一般に時間に依存する．微分同相変換のもとで不変でないラグランジアンを使うことに，読者は不安を感じるかもしれない．しかし，通常のゲージ理論と同様に，微分同相不変性は明らかに回復できる．パラメータ$\xi^0(t,\vec x)$を持つ時間微分同相変換を行い，そのパラメータを，時間微分同相変換のもとで$\pi(t,\vec{x})\rightarrow \pi(t,\vec{x})-\xi^0(t,\vec{x})$とシフトする場$\pi(t, \vec x)$へ昇格させればよい．スカラー$\pi$は時間微分同相変換を非線形に実現するゴールドストーン・モードであり，FRW解のまわりのスカラー摂動を記述する．

質量を持つゲージ・ボソンの縦成分の物理は，十分高いエネルギーでは，スカラーのゴールドストーン・モードに着目して調べられることがよく知られている（等価定理）．ここでも同じことが成り立つ．十分高いエネルギーでは重力との混合は重要でなくなり，ゴールドストーン・モードに集中できる．この領域では物理が非常に見通しよくなり，宇宙論的摂動についての情報の大部分を得ることができる．ラグランジアン(\ref{eq:Laguni})に対して破れた微分同相変換を行い，ゴールドストーン・モード$\pi$に着目すると，次を得る：
\begin{eqnarray} \label{theeffectiveaction} &&S_\pi = \int \! d^4 x \:
\sqrt{-g}\left[ M_{\rm Pl}^2 \dot H
 \, (\partial_\mu \pi)^2\right.\nonumber \\ &&\qquad\left.+2 M^4_2\left(\dot\pi^2+ \dot\pi^3-\dot\pi \frac{1}{a^2}(\di_i\pi)^2\right)-\frac{ 4}{3} M^4_3 \dot\pi^3 - \frac{\bar M ^2}{2} \,
\frac{1}{a^4}(\di_i ^2 \pi)^2 +\ldots
\right].
 \end{eqnarray}

ユニタリーゲージにおけるそれぞれの不変演算子は，ゴールドストーン場の（非線形な）演算子へと昇格する．微分同相不変性の非線形実現によって，各種の項のあいだに関係が強制されるのである．

理論の系統的な構成に進む前に，このアプローチの利点を簡単に述べておこう．
\begin{itemize}
\item 標準的な運動項とスローロール・ポテンシャルを持つスカラー場による，最も標準的なインフレーションのシナリオから出発し，ありうるすべての高エネルギー効果についてのわれわれの無知を，主要な不変演算子によってパラメータ化した．たとえば摂動の非ガウス性の測定や，重力波のスペクトルの傾きに対する整合性関係からのずれを調べることによって，実験は各演算子に制限を与える．ある意味で，これは素粒子物理学で行うことと類似している．そこでは，標準模型からのずれを記述し，したがって新しい物理の効果を表す演算子の大きさに制限を課す．

\item 対称性と，背景がインフレーションを起こすという要請によって何が強制され，何が自由であるかが明確になる．たとえば式(\ref{theeffectiveaction})は，空間運動項$(\nabla \pi)^2$が$\dot H$に比例する一方，時間運動項$\dot\pi^2$は自由であることを示している．もう一つの例は，ユニタリーゲージの演算子$(g^{00}+1)^2$である．これをゴールドストーン場$\pi$で書くと，$\pi$励起の音速を低下させる2次項$\dot\pi^2$と，モード間の相互作用，すなわち非ガウス性を増大させる3次項$\dot\pi (\nabla\pi)^2$が生じる．したがって，他の演算子との相殺がなければ，音速の低下と非ガウス性の増大は対称性によって結び付けられる．さらに，$\pi$励起が光円錐の外へ伝播することを避けるため，この演算子の係数は正でなければならない．

\item 可能な演算子をすべて把握できる．たとえば微分の最低次では，3つの$\pi$モード間の相互作用は$(g^{00}+1)^2$と$(g^{00}+1)^3$によって変更できる．これは3点関数の異なる2つの形状に対応し，原理的には実験によって区別して，各演算子の大きさを決めることができる．

\item 可能な単一場模型がすべて統一される．たとえば，かなり特徴的な予言を持つDBIインフレーション\cite{Alishahiha:2004eh,Chen:2005ad,Shandera:2006ax,Kecskemeti:2006cg,Shiu:2006kj}のように，ラグランジアンを$L((\partial\phi)^2,\phi)$へ変更した模型が関心を集めてきた．われわれの記述では，これらは演算子$(g^{00}+1)^n$が大きい場合に対応する．もう一つの興味深い極限は$\dot H \to 0$である．この場合，主要な空間運動項は$\bar M^2$に比例する演算子から生じ，$(\nabla^2 \pi)^2$という形になる．この極限はゴースト・インフレーション\cite{Arkani-Hamed:2003uz}を記述する．

\item $\phi$による記述では，場の再定義$\phi \to \tilde\phi(\phi)$を行うことができる．得られるラグランジアンが同じ物理を記述することは確かだが，そのことは自明ではない．簡単な例は，ラグランジアン
\be
f(\phi)^2 (\partial\phi)^2 -V(\phi) \;,
\ee
で与えられる．ここで$f$は一般の関数である．これは，場の再定義$\tilde\phi(\phi)$，$d\tilde\phi/d\phi = f(\phi)$を通じて，標準的な運動項と別のポテンシャルを持つラグランジアンに等価である．しかし，より一般的な項を考えると，異なるラグランジアン間の等価性はさらに複雑になる．一方，時間微分同相変換の一種の標準的な非線形表現を実現する$\pi$のレベルには，この曖昧さは存在しない．

\item $\phi$による記述では，演算子に含まれる場の一部を背景解上で評価できるため，摂動を調べる際に演算子の重要性をどう評価すべきかは明らかでない．たとえば$(\partial\phi)^{2n}$という形のすべての演算子を含む理論では，DBIインフレーションで起こるように，背景の速度$\dot\phi_0$が十分大きければ，それらがすべて同じ重要性を持ちうる．一方，$\pi$による記述では，演算子の重要性を見積もる通常の方法が機能する．このことはゴースト・インフレーションではさらに明瞭であり，$\pi$が非相対論的な分散関係を持つため，演算子のスケーリングは$\pi$による記述ではじめて明確になる．

\item 宇宙論的摂動のループ補正を計算する際には，解のまわりで演算子を直接パラメータ化することが決定的に重要である．外線数が与えられたダイアグラムは，一般にUV発散を含む．ラグランジアン(\ref{eq:Laguni})では，$n$個の摂動間の相互作用を記述する項が有限個しかないので，この発散は容易に繰り込める．一方，$\phi$のラグランジアンのレベルでは，$n$個の摂動間の相互作用に寄与する演算子が\emph{無限個}存在する．実際，各演算子の多数の脚を背景上に置くことができるため，演算子と摂動のダイアグラムとの関係はかなり不明瞭になる．
\end{itemize}

\sectionlink{ユニタリーゲージにおける作用の構成}
インフレーションとは，ハッブル・パラメータがほぼ一定である宇宙の加速膨張期である．この準デ・ジッター背景には，物理的な時計によって定まる特別な空間的スライシングがあり，これによって減速膨張する高温ビッグバンの発展へ滑らかに接続できる．このスライシングは通常，時間発展するスカラー$\phi(t)$によって実現される．念頭に置いておくとよいもう一つの例は完全流体である\footnote{実際，たとえば文献\cite{Dubovsky:2005xd}で示されているように，完全流体の渦なしの励起は，微分結合するスカラーによって記述できる．}．この解のまわりの摂動を記述するため，特別なスライシングが$t$一定面と一致するゲージ，すなわち$\delta\phi(\vec x,t)=0$を選ぶことができる．このゲージでは明示的なスカラー摂動はなく，計量の揺らぎだけが存在する．時間微分同相変換は固定され，もはやゲージ対称性ではないので，重力子は3つの自由度を記述する．スカラー摂動は計量に吸収されたのである．

このゲージにおける最も一般的なラグランジアンは何だろうか．計量$g_{\mu\nu}$の関数であり，（線形に実現された）時間依存する空間微分同相変換$x^i\rightarrow x^i+\xi^i(t,\vec{x})$のもとで不変な演算子を書き下さなければならない．実際，空間微分同相変換は破れていない．すべての微分同相変換のもとで不変なリーマン・テンソルを含む通常の項に加え，系の対称性が小さくなっているため，さらに多くの項が許される．それらは，重力子に吸収された追加の自由度を記述する．たとえば，$g^{00}$が空間微分同相変換のもとでスカラーであり，ユニタリーゲージのラグランジアンに自由に現れうることは容易に分かる．微分を含まない項は，$g^{00}$の多項式に限られる．時空に特別なスライシングがあるので，そのスライシングを記述する幾何学的な量を書くことも許される．たとえば，時間一定面の外的曲率$K_{\mu\nu}$は空間微分同相変換のもとでテンソルであり，作用中に用いることができる．作用のどの項にも，一般の時間の関数を掛けられることに注意しよう．\ref{app:generic}では，最も一般的なラグランジアンが次のように書けることを証明する：
\begin{eqnarray}
\label{eq:action}\nonumber
S & = & \int  \! d^4 x \; \sqrt{- g} \Big[ \frac12 M_{\rm Pl}^2 R - c(t)
g^{00} -  \Lambda(t) + \frac{1}{2!}M_2(t)^4(g^{00}+1)^2+\frac{1}{3!}M_3(t)^4 (g^{00}+1)^3+ \\
&& - \frac{\bar M_1(t)^3}{2} (g^{00}+1)\delta K^\mu {}_\mu
-\frac{\bar M_2(t)^2}{2} \delta K^\mu {}_\mu {}^2
-\frac{\bar M_3(t)^2}{2} \delta K^\mu {}_\nu \delta K^\nu {}_\mu + ...
\Big] \; ,
\end{eqnarray}
ここで省略記号は，揺らぎについて高次の項，またはより多くの微分を持つ項を表す．非摂動FRW背景からの時間一定面の外的曲率の変化を$\delta K_{\mu\nu}$と書く．すなわち，$\delta K_{\mu\nu} = K_{\mu\nu} - a^2 H h_{\mu\nu}$であり，$h_{\mu\nu}$は誘導空間計量である．上の作用で，選んだFRW解のまわりの1次の摂動を含むのは最初の3項だけであり，他の項はすべて明示的に2次以上であることに注意しよう．したがって，係数$c(t)$と$\Lambda(t)$は，指定されたFRWの発展$H(t)$を持つという要請，すなわちこの解のまわりでタドポール項が相殺するという要請によって固定される．これらの係数を固定する前に，この単純化が自明ではないことを理解しておくことが重要である．背景解のまわりで1次の寄与を与える演算子は，無限個あると予想されるかもしれない．しかし作用は，$\delta K_{\mu\nu}$や$g^{00}+1$のような1次の項の多項式として書けるので，演算子が1次，2次，あるいはそれ以上の次数から始まるかが明確になる．式(\ref{eq:action})の項以外のすべての1次項は微分を含み，部分積分によって，ここで考えた3つの1次項と高次の共変な項との組合せにできる．この構成は\ref{app:tadpoles}で具体的に行う．したがって，\emph{非摂動の膨張史によって$c(t)$と$\Lambda(t)$が固定され，異なる模型間の違いは高次の項に表される}．

指定されたFRWの発展が解であることを課して，1次項を固定しよう．すでに述べたように，$c$と$\Lambda$に比例する項だけが，摂動の0次で消えず，したがってアインシュタイン方程式の右辺に寄与するエネルギー運動量テンソル
\begin{equation}
T_{\mu \nu} = -\frac{2}{\sqrt{-g}}\frac{\delta S_{\rm
matter}}{\delta g^{\mu\nu}}
\end{equation}
を与える．インフレーション中に主として関心があるのは，平坦なFRW宇宙である（一般の場合については\ref{app:tadpoles}を参照）：
\be
ds^2 = -dt^2 + a^2(t) d \vec{x}^2 
\ee
このため，フリードマン方程式は次で与えられる：
\begin{eqnarray}
H^2 & = & \frac{1}{3 M_{\rm Pl}^2} \big[ c(t)+\Lambda(t)\big]  \\
\frac{\ddot a}{a} = \dot H + H^2 & =  & -\frac{1}{3 M_{\rm Pl}^2} \big[ 2
c(t)-\Lambda(t) \big] \;.
\end{eqnarray}
$c$と$\Lambda$について解くと，作用(\ref{eq:action})を次のように書き換えられる：
\begin{eqnarray}
\label{eq:actiontad}
S &=& \int  \! d^4 x \; \sqrt{- g} \Big[ \frac12 M_{\rm Pl}^2 R + M_{\rm Pl}^2 \dot H
g^{00} - M_{\rm Pl}^2 (3 H^2 + \dot H) +  \nonumber \\ && \frac{1}{2!}M_2(t)^4(g^{00}+1)^2+\frac{1}{3!}M_3(t)^4 (g^{00}+1)^3+ \nonumber \\
&& - \frac{\bar M_1(t)^3}{2} (g^{00}+1)\delta K^\mu {}_\mu
-\frac{\bar M_2(t)^2}{2} \delta K^\mu {}_\mu {}^2
-\frac{\bar M_3(t)^2}{2} \delta K^\mu {}_\nu \delta K^\nu {}_\mu + ...
\Big] \; .
\end{eqnarray}
すでに述べたように，上の作用に含まれる演算子の係数はすべて，一般の時間依存性を持ちうる．しかし，ここで関心があるのは，1ハッブル時間のあいだに$H$と$\dot H$が大きく変化しない解である．したがって，他のすべての演算子についても同じことが成り立つと仮定するのが自然である．この仮定のもとで，ラグランジアンは近似的に時間並進不変になる\footnote{時間シフトが厳密な対称性となる極限では$\dot H \to 0$となるため，注意して極限を取らなければならない．ゴースト・インフレーションの議論で見るように，これはゴールドストーン場の空間運動項が消えることを意味する．}．したがって，ループ効果によって生じる時間依存性は，小さな対称性の破れのパラメータによって抑制される\footnote{この対称性は，時間微分同相変換の破れとは無関係であることに注意しよう．$\phi$による記述でこの対称性がどう現れるかを見るには，適切な場の再定義を行えば，常に$\dot\phi =$ const.と仮定できることに注目すればよい．この選択のもとでは，シフト対称性$\phi \to \phi $ + const.から，ユニタリーゲージのラグランジアンの時間並進不変性が従う．この対称性と$\phi$のラグランジアンの時間並進対称性は，背景によって対角部分群へと破られる．この残存対称性が，ユニタリーゲージのラグランジアンにおける時間シフトである．}．この仮定はとくに都合がよい．係数が急速に時間変化すると，指数関数的膨張が生む摩擦に打ち勝つことがあり，その結果，標準的なFRW宇宙論の一様性問題を解決するために必要な力学的アトラクターとしての性質を，インフレーションが失うおそれがあるからである．

このアプローチが記述するのは，スカラー・モードだけでなく重力についても最も一般的なラグランジアンであることを強調しておく．たとえば，高エネルギー効果はリーマン・テンソルの摂動$\delta R_{\mu\nu\rho\sigma}$を含む演算子に表される．これらの補正は微分について高次であるため，以下では明示的には扱わない．

単純なインフレーション模型をこの記述でどのように書くか，いくつか例を挙げよう．標準的な運動項とスローロール・ポテンシャル$V(\phi)$を持つ模型は，ユニタリーゲージで次のように書ける：
\be
\int \! d^4x \: \sqrt{-g}
\left[ -\frac 1 2 (\partial \phi)^2 - V(\phi) \right]  \to  \int \!
d^4x \: \sqrt{- g} \left[ -\frac{\dot \phi_0(t)^2}{2} g^{00} - V(\phi_0(t))
\right] \; .
\ee
フリードマン方程式から$\dot\phi_0(t)^2=-2M^2_P \dot{H}$および$V(\phi(t))=M_{\rm Pl}^2 (3H^2+\dot H$)が得られるので，作用は，最初の3項以外をすべてゼロとした式(\ref{eq:actiontad})の形であることが分かる．ただし，他のすべての項はループ補正によって生成されるため，これが厳密に成り立つことは明らかにない．それらの項は，この単純なスローロール・インフレーション模型に対する高エネルギー物理の，ありうるすべての効果を表している．

さらに一般的な場合として，各$\phi$に作用する微分が高々1つであるすべてのラグランジアンを含めることができる．すなわち，$L= P(X,\phi)$，ただし$X=g^{\mu\nu}\partial_\mu\phi\partial_\nu\phi$である．非摂動解$\phi_0(t)$のまわりでは，
\be
S = \int \! d^4x \: \sqrt{-g} \;  P(\dot\phi_0(t)^2 g^{00}, \phi(t))
\ee
となる．これは明らかに先ほどの形であり，$M_n^4(t) =\dot\phi_0(t)^{2n} \partial^n P/\partial X^n$を$\phi_0(t)$で評価すればよい．外的曲率を含む項では，1つのスカラーに2つ以上の微分が作用し，これらの項は厳密なデ・ジッター極限$\dot H \to 0$で決定的に重要になる．


\sectionlink{ゴールドストーン・ボソンの作用\label{sec:Goldstone}}
ユニタリーゲージのラグランジアンは，重力子の2つのヘリシティと1つのスカラー・モードという3つの自由度を記述する．破れた時間微分同相変換を行うと（シュテュッケルベルクの方法），このスカラー・モードは，その対称性を非線形に実現するゴールドストーン・ボソンとして明示的に現れる．質量を持つゲージ・ボソンの縦成分に対する等価定理\cite{Cornwall:1974km}との類推から，混合を無視できる短距離では，ゴールドストーン場の物理は重力子の2つのヘリシティから分離すると予想される．同じ方法をここで適用する前に，非アーベル・ゲージ理論で何が起こるかを簡単に振り返ろう．

非アーベル・ゲージ場$A_\mu^a$のユニタリーゲージにおける作用は
\be
S = \int \! d^4x  \,-\frac{1}{4} {\rm Tr}\, F_{\mu\nu}F^{\mu\nu}-\frac12 m^2 {\rm Tr}\, A_\mu A^\mu \ ,
\ee
であり，ここで$A_\mu = A_\mu^a T^a$である．ゲージ変換のもとで
\be
\label{eq:AmuU}
A_\mu \to U A_\mu U^\dagger + \frac{i}{g} U \partial_\mu U^\dagger \equiv \frac{i}{g} U D_\mu U^\dagger \;.
\ee
となる．したがって作用は
\be
S = \int \! d^4x  \,-\frac{1}{4} {\rm Tr}\, F_{\mu\nu}F^{\mu\nu} - \frac12 \frac{m^2}{g^2} {\rm Tr} D_\mu U^\dagger D_\mu U \;.
\ee
となる．$U=\exp{[i T^a \pi^a(t,\vec x)]}$と書くことによって，ゲージ不変性を「回復」できる．ここで$\pi^a$はスカラー（ゴールドストーン場）であり，ゲージ変換$\Lambda$のもとで次のように非線形に変換する：
\be
e^{i T^a \widetilde\pi^a(t, \vec x)} = \Lambda(t, \vec x) \,e^{i T^a \pi^a(t,\vec x)}
\ee
正準規格化$\pi_c \equiv m/g \cdot \pi$へ移ると，ゴールドストーン・ボソンの自己相互作用はスケール$4 \pi m/g$で強結合になることが分かる．このスケールは，パラメータ依存性の意味でゲージ・ボソンの質量より高い．ゴールドストーン場を再導入する利点は，エネルギー$E \gg m$で，ゴールドストーン場とゲージ場の横成分との混合が重要でなくなり，2つのセクターが分離することである．実際，式(\ref{eq:AmuU})の混合項は
\be
\frac{m^2}{g} A_{\mu}^a \partial^\mu \pi^a = m A_{\mu}^a \partial^\mu \pi_c^a
\ee
という形を取り，$E \gg m$では正準運動項$(\partial \pi_c)^2$に比べて重要でなくなる．領域$m \ll E \ll 4 \pi m /g$では，ゴールドストーン場$\pi$の物理は弱結合であり，横成分との混合を無視して調べることができる．

時間微分同相変換が破れた場合について，同じ手順をたどろう．たとえば，次の2つの演算子に着目する：
\begin{equation}
\int d^4x\;  \sqrt{-g} \left[A(t)+B(t)g^{00}(x)\right] \ .
\end{equation}
破れた時間微分同相変換$t \to \widetilde t= t + \xi^0(x)$，$\vec{x} \to \vec{\widetilde{x}}=\vec{x}$のもとで，$g^{00}$は次のように変換する：
\begin{equation}
g^{00}(x)\to \widetilde g^{00}(\widetilde x(x))=\frac{\partial \widetilde x^0(x)}{\partial x^\mu}\frac{\partial \widetilde x^0(x)}{\partial x^\nu} g^{\mu\nu}(x) \,  .
\end{equation}
変換後の場で書いた作用は次で与えられる：
\begin{eqnarray}
\int d^4x\;  \sqrt{-\widetilde g(\widetilde x(x))} \left|\frac{\partial \widetilde x}{\partial x} \right| \left[A(t)+B(t) \frac{\partial x^0}{\partial \widetilde x^\mu}\frac{\partial x^0}{\partial \widetilde x^\nu} \widetilde g^{\mu\nu}(\widetilde x(x))\right]\ .
\end{eqnarray}
積分変数を$\widetilde x$へ変更すると，
\begin{eqnarray}
&&\int d^4\widetilde x\;  \sqrt{-\widetilde g(\widetilde x)} \left[A(\widetilde t-\xi^0(x(\widetilde x)))\right.\nonumber \\ &&\qquad\left.+B(\widetilde t-\xi^0(x(\widetilde x))) \frac{\partial (\widetilde t-\xi^0(x(\widetilde x)))}{\partial \widetilde x^\mu}\frac{\partial (\widetilde t-\xi^0(x(\widetilde x)))}{\partial \widetilde x^\nu} \widetilde g^{\mu\nu}(\widetilde x)\right].
\end{eqnarray}
を得る．ゴールドストーン場を再導入する手続きは，ゲージ理論の場合と同様である．上の作用に現れる$\xi^0$をすべて，
\begin{equation} 
\xi^0(x(\widetilde x)) \to - \widetilde \pi(\widetilde x ) \, .
\end{equation}
と置き換える．簡単のためチルダを省略すると，次を得る：
\begin{eqnarray}
\int d^4x\;  \sqrt{- g(x)} \left[A( t+\pi(x))+B(t+\pi(x)) \frac{\partial (t+\pi(x))}{\partial x^\mu}\frac{\partial (t+\pi(x))}{\partial x^\nu} g^{\mu\nu}(x)\right].
\end{eqnarray}
$\pi$に変換則
\begin{equation}
\pi(x) \to \widetilde\pi(\widetilde x(x))=\pi(x)-\xi^0(x) \ .
\end{equation}
を割り当てれば，上の作用が無限小変換だけでなく，すべての次数で微分同相変換のもとで不変となることを確認できる．この定義のもとで，$\pi$はスカラー場として変換するとともに，時間微分同相変換に対して追加のシフトを受ける．

この手続きをユニタリーゲージの作用(\ref{eq:actiontad})に適用すると，
\begin{eqnarray}\label{Smixed}
S = \int \! d^4 x  \: \sqrt{- g} &&\left[\frac{1}{2}M_{\rm Pl}^2 R
- M^2_{\rm Pl} \left(3H^2(t+\pi) +\dot{H}(t+\pi)\right)+ \right.\\ \nn
&&+M^2_{\rm
Pl} \dot{H}(t+\pi)\left(
(1+\dot\pi)^2g^{00}+2(1+\dot\pi)\partial_i\pi g^{0i}+
g^{ij}\partial_i\pi\partial_j\pi\right) + \\ \nonumber
&&\frac{M_2(t+\pi)^4}{2!}\left(
(1+\dot\pi)^2g^{00}+2(1+\dot\pi)\partial_i\pi g^{0i}+
g^{ij}\partial_i\pi\partial_j\pi+1\right)^2 + \nonumber\\
\nonumber && \left. \frac{M_3(t+\pi)^4}{3!}\left(
(1+\dot\pi)^2g^{00}+2(1+\dot\pi)\partial_i\pi g^{0i}+
g^{ij}\partial_i\pi\partial_j\pi+1\right)^3+ ... \right] \; ,
\end{eqnarray}
を得る．ここではひとまず，簡単のため外的曲率を含む項を無視した．

この作用はかなり複雑であり，この時点では，ユニタリーゲージのラグランジアンからゴールドストーン場$\pi$を再導入する利点は明らかでない．ゲージ理論の場合と同様に，十分短い距離では計量の揺らぎを無視してゴールドストーン場の物理を調べられるため，単純化が起こる．それが可能な領域は，ゲージ理論の場合と同じく，上のラグランジアンの混合項を見るだけで見積もれる．実際，式(\ref{Smixed})では，$\pi$と$g_{\mu\nu}$を混合する2次項の微分の数は$\pi$の運動項より少なく，ある高エネルギー・スケールより上で無視できることが分かる．一般に，その答えはどの演算子が存在するかに依存する．まず，タドポール項だけが重要な最も簡単な場合（$M_2=M_3=\ldots=0$）から始めよう．これは通常のスローロール・インフレーションに対応する．重力との主要な混合は，次の形の項から生じる：
\begin{equation}
\sim M_{\rm Pl}^2 \dot H \dot\pi \delta g^{00} \ .
\end{equation}

正準規格化（$\pi_c\sim M_{\rm Pl} \dot H^{1/2}\pi,\  \delta g_c^{00}\sim M_{\rm Pl} \delta g^{00}$）を行うと，$E_{\rm mix}\sim \epsilon^{1/2} H$より高いエネルギーでは混合項を無視できることが分かる．ここで$\epsilon$は通常のスローロール・パラメータ$\epsilon\equiv -\dot H/H^2$である．もう一つ関心があるのは，$M_2$の演算子が大きくなる場合である．この場合には，
\begin{equation}
\sim M_2^4 \dot \pi \delta g^{00}
\end{equation}
という形の混合項があり，正準規格化を行うと（今度は$\pi_c\sim M_2^2\pi$であることに注意），$E_{\rm mix}\sim M_2^2/M_{\rm Pl}$より高いエネルギーで無視できるようになる\footnote{ここで調べている理論ではローレンツ対称性が自発的に破れているので，混合を無視できるエネルギーと運動量の領域をそれぞれ別に定義すべきである．宇宙論的摂動については，エネルギーの範囲だけに関心を持つ．}．

領域$E\gg E_{\rm mix}$では，作用は大幅に単純化され，次のようになる：
\begin{eqnarray}\label{Spi} 
&&\! \! S_{\rm \pi} = \int \!  d^4 x   \sqrt{- g} \left[\frac12 M_{\rm Pl}^2 R -M^2_{\rm Pl}
\dot{H} \left(\dot\pi^2-\frac{ (\partial_i \pi)^2}{a^2}\right)
\right.\nonumber \\ &&\qquad\left.+2 M^4_2
\left(\dot\pi^2+\dot{\pi}^3-\dot\pi\frac{(\partial_i\pi)^2}{a^2}
\right) -\frac{4}{3} M^4_3 \dot{\pi}^3+ ... \right].
\end{eqnarray}

インフレーション模型が与えられたとき，現在の宇宙論的観測に対する予言を計算したい．この観点からは，デカップリング極限(\ref{Spi})は，このような極端な赤外スケールにはまったく無関係に思える．しかし，通常の単一場スローロール・インフレーションの場合と同じく，摂動論の任意の次数で地平線の外側において一定となる量，通常の$\zeta$変数が存在することを証明できる\cite{Salopek:1990jq,Lyth:2004gb}（高階の空間微分を持つ項を含めた一般化については，文献\cite{verification}の付録Dを参照）．保存量が存在する直観的な理由は，地平線を出た後には，異なる領域がまったく同じように発展することにある．唯一の違いは，ある領域が別の領域に比べてどれだけ膨張したかという点であり，この差が一定に保たれる．

したがって，問題は地平線通過直後の相関関数を計算することに帰着する．そこで，$H$程度の赤外（IR）エネルギー・カットオフを持つラグランジアンを調べればよい．デカップリング・スケール$E_{\rm mix}$が$H$より小さければ，$\pi$のラグランジアン(\ref{Spi})は，$E_{\rm mix}/H$で抑制された項を除いて正しい予言を与える．

すでに述べたように，ユニタリーゲージのラグランジアンの係数の時間依存性は，ハッブル時間に比べて緩やかであり，何らかの一般化されたスローロール・パラメータによって抑制されると仮定している．これは，係数のテイラー展開から生じる追加の$\pi$の項が小さいことを意味する．とくに，3次項のような関連演算子，すなわち赤外へ向かうにつれて支配的になる演算子は，スケール$H$では重要でなく，そのためラグランジアン(\ref{Spi})では無視してある．

以上から，ラグランジアン(\ref{Spi})を用いれば，重力との混合に支配されないすべての観測量を計算できる．混合に支配される量の例には，通常のスローロール・インフレーションにおける非ガウス性がある\cite{Maldacena:2002vr,Seery:2006vu}．ただし，スペクトルの傾きは，主要次数ではラグランジアン(\ref{Spi})から計算できることに注意しよう．後で見るように，その値は実際，混合項を無視して計算した地平線通過時のパワー・スペクトルから単純に導ける．重力との混合が重要になり，ゴールドストーン場の作用だけでは予言に不十分となる場合にも，われわれのアプローチが妥当性を失うわけではないことを強調しておく．作用(\ref{eq:actiontad})は模型についてのすべての情報を含んでおり，重力との混合が大きい場合にも，これを用いてすべての予言を計算できる．

